% P108: MANNA, ADASS XXXVI (2026) poster proceedings.
% Built from the ADASS_template.tex LaTeX file, 24 Oct 2025, itself based on the
% ASP general author template (aspauthor.tex, Copyright 2021, Astronomical
% Society of the Pacific Conference Series, revision 29 November 2021).
% Condensed from the AAS manuscript in dangause/manna-paper (main, 19064da).
%
% Build and check from this folder:  make pdf  and  make check
% ADASS rules: no \usepackage other than asp2023, no \newcommand or
% \renewcommand, PNG/JPG figures named P108_f<n>, 4 pages including references.

\documentclass[11pt,twoside]{article}

% Do NOT use ANY packages other than asp2023.
\usepackage{asp2023}

\aspSuppressVolSlug
\resetcounters

\bibliographystyle{asp2021}

\markboth{Gause et al.}{MANNA: An MCP Architecture for Astronomical Archives}

\begin{document}

\title{MANNA: An MCP Architecture for NOIRLab, NRAO, and Additional Archives}

\author{Daniel~Gause,$^1$ Robert~Nikutta,$^2$ Brian~Mason,$^1$ Adele~Plunkett,$^1$ Ryan~Loomis,$^1$ Stephanie~Juneau,$^2$ Stella~S.~R.~Offner,$^3$ and Eric~J.~Murphy$^1$}
\affil{$^1$National Radio Astronomy Observatory, Charlottesville, VA, USA; \email{dgause@nrao.edu}}
\affil{$^2$NSF NOIRLab, Tucson, AZ, USA}
\affil{$^3$Department of Astronomy, The University of Texas at Austin, Austin, TX, USA}

% This section is for ADS Processing. One line per author; paperauthor has 9 arguments.
\paperauthor{Daniel~Gause}{dgause@nrao.edu}{0009-0004-8062-3810}{National Radio Astronomy Observatory}{}{Charlottesville}{VA}{22903}{USA}
\paperauthor{Robert~Nikutta}{robert.nikutta@noirlab.edu}{0000-0002-7052-6900}{NSF NOIRLab}{}{Tucson}{AZ}{85719}{USA}
\paperauthor{Brian~Mason}{bmason@nrao.edu}{0000-0002-8472-836X}{National Radio Astronomy Observatory}{}{Charlottesville}{VA}{22903}{USA}
\paperauthor{Adele~Plunkett}{aplunket@nrao.edu}{0000-0002-9912-5705}{National Radio Astronomy Observatory}{}{Charlottesville}{VA}{22903}{USA}
\paperauthor{Ryan~Loomis}{rloomis@nrao.edu}{0000-0002-8932-1219}{National Radio Astronomy Observatory}{}{Charlottesville}{VA}{22903}{USA}
\paperauthor{Stephanie~Juneau}{stephanie.juneau@noirlab.edu}{0000-0002-0000-2394}{NSF NOIRLab}{}{Tucson}{AZ}{85719}{USA}
\paperauthor{Stella~S.~R.~Offner}{soffner@utexas.edu}{0000-0003-1252-9916}{The University of Texas at Austin}{Department of Astronomy}{Austin}{TX}{78712}{USA}
\paperauthor{Eric~J.~Murphy}{emurphy@nrao.edu}{0000-0001-7089-7325}{National Radio Astronomy Observatory}{}{Charlottesville}{VA}{22903}{USA}

% Author index entries, surname first.
%\aindex{Gause,~D.}
%\aindex{Nikutta,~R.}
%\aindex{Mason,~B.}
%\aindex{Plunkett,~A.}
%\aindex{Loomis,~R.}
%\aindex{Juneau,~S.}
%\aindex{Offner,~S.~S.~R.}
%\aindex{Murphy,~E.~J.}

% Subject index entries, all existing terms from Index.py.
%\ssindex{International Virtual Observatory Alliance (IVOA)}
%\ssindex{databases!query language!ADQL}
%\ssindex{archives!individual!ALMA Science Archive}
%\ssindex{archives!multiple!National Radio Astronomy Observatory (NRAO) data archive}
%\ssindex{Artificial intelligence}

\begin{abstract}
While many astronomical data archives implement the same set of standards defined by the International Virtual Observatory Alliance (IVOA), differences in underlying data and architectural implementations necessitate significant archive-specific knowledge. We present MANNA (MCP Architecture for NOIRLab, NRAO, and Additional Archives), a Model Context Protocol (MCP) server that provides AI agents with tools for archive discovery and data querying. With a custom evaluation framework, we ran archive tasks with and without MANNA for one open-weight and three Anthropic Claude models. MANNA reduced the time needed per solved task for all four models and substantially increased the number of tasks passed by the smallest Claude model. The differences were larger when the time and turns per task were limited. MANNA is compatible with all IVOA-based astronomical archives.
\end{abstract}

\section{Introduction}
\label{sec:intro}
Most large astronomical archives offer programmatic access through services defined by the International Virtual Observatory Alliance (IVOA), such as the Table Access Protocol (TAP). In practice these standards have been implemented unevenly: archives differ in the tables they expose, the column names they use, the ADQL functions they support, and whether a query must run asynchronously. Using an archive through its IVOA services therefore still requires knowledge of that particular archive.

AI agents built on large language models can write and submit such queries, but what they know about an archive comes from their training data and from public documentation, both of which may be incomplete or out of date. An agent can often find the coordinates of a target or the name of the archive that holds its data, yet fail to retrieve the data itself, because it lacks details such as the address of the query service, the table that holds the observations, or the query syntax the service accepts. This necessitates expensive iteration through trial and error, and can fundamentally limit agentic capabilities. This paper presents MANNA (MCP Architecture for NOIRLab, NRAO, and Additional Archives), an MCP server designed to help AI agents overcome these limitations. The Model Context Protocol \citep[MCP;][]{Anthropic2024mcp} is an open standard for connecting AI agents to external tools and data sources. Section~\ref{sec:arch} describes the server, Section~\ref{sec:eval} its evaluation with four language models, and Section~\ref{sec:discussion} our conclusions.

\section{Architecture}
\label{sec:arch}
MANNA is a Python MCP server built with the FastMCP framework \citep{fastmcp}. It consists of tools (core and workflow) designed to facilitate the use of astronomical archives' IVOA services, security- and context-window-aware result handling, and archive notes that provide supplementary information on specific archives (Figure~\ref{fig:arch}).

\articlefigure[width=\textwidth]{P108_f1.png}{fig:arch}{Schematic of MANNA. The AI agent calls MANNA's tools over MCP; the archive notes supply archive-specific knowledge to the tools and are re-checked against the archives by the note audit.}

{\emergencystretch=1.5em MANNA's core tools cover each step of retrieving data through IVOA services, from resolving a target name to following an asynchronous job, and include \texttt{describe\_table}, \texttt{resolve\_target\_name}, \texttt{search\_ivoa\_registry}, and \texttt{run\_adql\_query}, which submits queries to TAP services \citep{2019ivoa.spec.0927D}. Its workflow tools chain several of these steps into one call, which reduces the number of decisions the agent has to make; \texttt{count\_\allowbreak observations\_\allowbreak near\_target}, for example, resolves a target, selects an archive, and runs that archive's positional count.\par}

By default, \texttt{run\_adql\_query} submits an ADQL query \citep{2023ivoa.spec.1215M} synchronously. If the archive has not answered within 20 seconds, or the result exceeds 200 rows or 48\,KB, MANNA resubmits the query as an asynchronous job and returns the job URL with a short Python recipe for fetching the result, so that large tables stay out of the agent's context window.

The archive notes build on the IVOA discovery mechanisms (the registry, service capabilities, and \texttt{TAP\_SCHEMA}, which the tools still query at run time) and add what these do not express. ALMA tables, for instance, list each dataset once per spectral window, so a count of observations must count distinct data\-sets; Astro Data Lab requires its own \texttt{q3c} functions for positional queries. Each note carries a check: an ADQL query with an expected outcome, a test that named columns still exist, or, for the 36 of 74 notes that no single query can test, a flag for review by hand. A weekly note audit runs the automated checks against the live archives and opens an issue for any note that no longer holds.

MANNA includes notes for seven archives from radio to optical wavelengths, most fully for NOIRLab's Astro Data Lab \citep{2020A&C....3300411N}, the ALMA Science Archive \citep{2025ASPC..538..206S}, and the NRAO Science Data Archive, which at NRAO's request MANNA has not queried by default since September 2026, while its TAP service is rebuilt.

\section{Evaluation and Results}
\label{sec:eval}
We measured MANNA's effect with a custom evaluation framework, distributed with the code, that runs Claude Code \citep{claudecode} in headless mode with the model under test and records every tool call. Each task pairs a request with a grading specification: the tools that must appear among the calls, constraints on their arguments, and either a ground-truth answer or, for open-ended questions, a rubric scored by a judge model.

The tool-use evaluation (Program I) tests whether a model uses MANNA's tools correctly. The value-added measurement (Program II) runs a separate set of 26 tasks, such as counting the distinct ALMA observations of NGC 253, with MANNA and without it; without MANNA the agent keeps Claude Code's built-in tools for shell commands and web access.

We ran both programs on MANNA v0.10.0 through Claude Code in September 2026 with Anthropic's Haiku 4.5 \citep{Anthropic2025haiku45}, Sonnet 5 \citep{Anthropic2026sonnet5}, and Opus 4.8 \citep{Anthropic2026opus48}, and NVIDIA's open-weight Nemotron 3 Super model \citep{2026arXiv260412374N}, which we hosted ourselves and also used as the judge. Sonnet 5 and Opus 4.8 ran at Claude Code's default effort level, \texttt{high}. Costs are Claude Code's estimates at list prices; the self-hosted Nemotron 3 Super has no comparable cost. Program II tasks were allowed 1 hour and 200 turns, and each model and condition was run three times; we also made one capped run of each, limited to 600 seconds and 20 turns per task. Program I was run once per model. Because the NRAO archive is paused, we ran only the 14 of 26 Program I tasks and 23 of 26 Program II tasks that do not need it. All results are archived on Zenodo \citep{manna2026evals}.

In Program I, Sonnet 5 and Opus 4.8 passed all 14 tasks and Haiku 4.5 and Nemotron 3 Super passed 12. Table~\ref{tab:p2} summarizes Program II. Without limits, every model except Haiku passed nearly all tasks with or without MANNA, and MANNA's effect was mainly on the work needed per solved task: without it, Nemotron made 15.5 times as many tool calls and took 8.1 times as long. For Haiku, MANNA also changed the outcome: it passed 21.7 of 23 tasks on average with MANNA and 15.0 without.

% Values: n3-summary.md (2026-09-22-subset-evals), uncapped means, per-solved-task
% multipliers, and the 600 s / 20-turn capped table; also AAS Tables 4 and 6.
\begin{table}[!ht]
\caption{Program II through Claude Code. Tasks passed (of 23) are means of three uncapped runs (1\,h and 200 turns per task) or one capped run (600\,s and 20 turns). The last three columns are ratios per solved task, without MANNA to with it, in the uncapped runs; values above 1 favor MANNA.}
\label{tab:p2}
\smallskip
\begin{center}
{\small
\begin{tabular}{lccccccc}
\tableline
\noalign{\smallskip}
 & \multicolumn{2}{c}{Uncapped} & \multicolumn{2}{c}{Capped} & \multicolumn{3}{c}{Per solved task}\\
Model & With & Without & With & Without & Calls & Time & Cost\\
\noalign{\smallskip}
\tableline
\noalign{\smallskip}
Nemotron 3 Super & 21.0 & 20.3 & 22 & 12 & 15.5$\times$ & 8.1$\times$ & --\\
Haiku 4.5        & 21.7 & 15.0 & 22 & 14 & 3.1$\times$  & 2.3$\times$ & 6.7$\times$\\
Sonnet 5         & 22.0 & 21.3 & 23 & 18 & 1.1$\times$  & 2.2$\times$ & 1.3$\times$\\
Opus 4.8         & 23.0 & 23.0 & 23 & 20 & 2.3$\times$  & 3.4$\times$ & 2.4$\times$\\
\noalign{\smallskip}
\tableline
\end{tabular}
}
\end{center}
\end{table}

The difference was largest under the capped limits. Without MANNA the models passed 12 to 20 of the 23 tasks and reached a limit on 3 to 11 of them; with MANNA they passed 22 or 23 and reached a limit on none.

\section{Discussion and Conclusions}
\label{sec:discussion}
MANNA narrowed the differences between models: with it, all four passed 21.0 to 23.0 of 23 tasks; without it, 15.0 to 23.0. Haiku 4.5 with MANNA passed 21.7 tasks for \$0.98 in total; Opus 4.8 without it passed all 23 for \$8.84. A smaller model with MANNA is therefore a practical alternative to a larger model without it.

The evaluation has limitations: Program II covers basic uses of archives, which a capable model can often handle by querying the services directly; each condition was run three times and the capped runs once; and the judge was itself a model under test.

\emergencystretch=1.5em
MANNA is documented at \url{https://manna.readthedocs.io}; it is open source under the MIT license \citep{manna2026repo} and on PyPI as \texttt{manna-mcp}. We encourage both the local use of MANNA by astronomers and its integration into astronomical applications and platforms, and we welcome the integration of additional IVOA-based archives via developer contributions to the repository.
\par\emergencystretch=0pt

\acknowledgments This work was supported by the NSF-Simons AI Institute for Cosmic Origins (CosmicAI; NSF Cooperative Agreement 2421782 and Simons Foundation award MPS-AI-00010515). This research uses services or data provided by the Astro Data Lab, which is part of the Community Science and Data Center (CSDC) Program of NSF NOIRLab. NOIRLab is operated by AURA, Inc. under a cooperative agreement with the U.S. National Science Foundation. The National Radio Astronomy Observatory is a facility of the U.S. National Science Foundation operated under cooperative agreement by Associated Universities, Inc.

\bibliography{P108}  % For BibTex

% if we have space left, we might add a conference photograph here. Leave commented for now.
% \bookpartphoto[width=1.0\textwidth]{foobar.jpg}{FooBar Photo (Photo: Any Photographer)}

\end{document}
